% !PRINTED 124
\weberpage{124}
%<<CONT

\gesperrt{so that here also the lower bound is a minimum.}

The \gesperrt{lower bound} of $X$ in any part of the region is either equal to $g$ or greater than $g$.

A domain of magnitudes determined by the inequalities
\begin{equation}
 |x|\le R,\qquad -R\le y\le \alpha
\tag{7}
\end{equation}
is represented in our Figure 2 by the segment $(PQ\alpha Q')$, which we shall call the segment $(P,\alpha)$.

% Exact 300-dpi crop from the registered Volume-I scan, printed Fig. 2.
\weberfigureimage{weber-b1-fig02.png}{0.62}

We now first determine a value $\eta$ of $y$ by a cut $(\mathfrak A,\mathfrak B)$ as follows: a number $\alpha$ between $-R$ and $+R$ is admitted into $\mathfrak A$ if the \gesperrt{lower bound} of $X$ in the segment $(P,\alpha)$ is greater than $g$, and a value $\beta$ is admitted into $\mathfrak B$ if the \gesperrt{lower bound} of $X$ in the segment $(P,\beta)$ is equal to $g$. This cut $(\mathfrak A,\mathfrak B)$ defines a number $\eta$ with the property that the \gesperrt{lower bound} of $X$ in the region $(P,y)$ is greater than $g$ if $y<\eta$, and equal to $g$ if $y>\eta$.

Now, by \S\,31, VI,
\[
 |f(\eta+iz)|
\]
is a continuous function of $z$ in the interval
\begin{equation}
 -\sqrt{R^2-\eta^2}\le z\le \sqrt{R^2-\eta^2},
\tag{8}
\end{equation}
which is denoted in Figure 2 by $M,M'$.

Thus by 3. this function receives, for some value $\zeta$ of $z$ in this interval, a minimum value $\gamma$, so that, if
\[
 \xi=\eta+i\zeta
\]
is put,
\begin{equation}
 |f(\xi)|=\gamma.
\tag{9}
\end{equation}
